\documentclass{article}%
\usepackage{amsmath}%
\setcounter{MaxMatrixCols}{30}%
\usepackage{amsfonts}%
\usepackage{amssymb}%
\usepackage{graphicx}
%TCIDATA{OutputFilter=latex2.dll}
%TCIDATA{Version=5.50.0.2953}
%TCIDATA{CSTFile=40 LaTeX article.cst}
%TCIDATA{Created=Monday, May 18, 2026 19:31:33}
%TCIDATA{LastRevised=Wednesday, May 20, 2026 15:31:13}
%TCIDATA{<META NAME="GraphicsSave" CONTENT="32">}
%TCIDATA{<META NAME="SaveForMode" CONTENT="1">}
%TCIDATA{BibliographyScheme=Manual}
%TCIDATA{<META NAME="DocumentShell" CONTENT="Standard LaTeX\Blank - Standard LaTeX Article">}
%BeginMSIPreambleData
\providecommand{\U}[1]{\protect\rule{.1in}{.1in}}
%EndMSIPreambleData
\newtheorem{theorem}{Theorem}
\newtheorem{acknowledgement}[theorem]{Acknowledgement}
\newtheorem{algorithm}[theorem]{Algorithm}
\newtheorem{axiom}[theorem]{Axiom}
\newtheorem{case}[theorem]{Case}
\newtheorem{claim}[theorem]{Claim}
\newtheorem{conclusion}[theorem]{Conclusion}
\newtheorem{condition}[theorem]{Condition}
\newtheorem{conjecture}[theorem]{Conjecture}
\newtheorem{corollary}[theorem]{Corollary}
\newtheorem{criterion}[theorem]{Criterion}
\newtheorem{definition}[theorem]{Definition}
\newtheorem{example}[theorem]{Example}
\newtheorem{exercise}[theorem]{Exercise}
\newtheorem{lemma}[theorem]{Lemma}
\newtheorem{notation}[theorem]{Notation}
\newtheorem{problem}[theorem]{Problem}
\newtheorem{proposition}[theorem]{Proposition}
\newtheorem{remark}[theorem]{Remark}
\newtheorem{solution}[theorem]{Solution}
\newtheorem{summary}[theorem]{Summary}
\newenvironment{proof}[1][Proof]{\noindent\textbf{#1.} }{\ \rule{0.5em}{0.5em}}
\begin{document}
\emph{FUNCION\ DE\ TRANSFERENCIA.}%

\[
R_{1}i_{1}\left(  t\right)  +\frac{1}{C}\int\left[  i_{1}\left(  t\right)
-i_{2}\left(  t\right)  \right]  dt=V_{e}\left(  t\right)
\]
%

\[
L\frac{di_{2}\left(  t\right)  }{dt}+\left(  R_{2}+R_{3}\right)  i_{2}\left(
t\right)  =\frac{1}{C}\int\left[  i_{1}\left(  t\right)  -i_{2}\left(
t\right)  \right]  dt
\]
%

\[
F_{s}\left(  t\right)  =V_{R3}\left(  t\right)  =R_{3}i_{2}\left(  t\right)
\]


\bigskip Aplicando transformada de Laplace:

$\left(  1\right)  $%
\[
R_{1}i_{1}\left(  s\right)  +\frac{1}{C_{s}}\int\left[  i_{1}\left(  s\right)
-i_{2}\left(  s\right)  \right]  =V_{e}\left(  s\right)
\]
$\ \ \ \ \ \ \left(  2\right)  $%
\[
L_{s}I_{2}\left(  s\right)  +\left(  R_{2}+R_{3}\right)  i_{2}\left(
s\right)  =\frac{1}{C}\left[  i_{1}\left(  s\right)  -i_{2}\left(  s\right)
\right]
\]


Al despejar la segunda ecuacion se agrupa el lado izquierdo:%

\[
L_{s}+\left(  R_{2}+R_{3}\right)  i_{2}\left(  s\right)  =\frac{1}{C}\left[
i_{1}\left(  s\right)  -i_{2}\left(  s\right)  \right]
\]


Multiplicamos por $C_{s}:$%

\[
C_{s}[L_{s}\left(  R_{2}+R_{3}\right)  ]i_{2}\left(  s\right)  =i_{1}\left(
s\right)  -i_{2}\left(  s\right)
\]


Pasamos $i_{2}\left(  s\right)  $ al otro lado:%

\[
i_{1}\left(  s\right)  =[C_{s}(L_{s}+R_{2}+R_{3}]+si+1]I_{2}\left(  s\right)
\]


Sustituir ec. (2) en ec. (1):

\bigskip%

\[
R_{1}i_{1}\left(  s\right)  +[L_{s}+\left(  R_{2}+R_{3}\right)  ]i_{2}\left(
s\right)  =V_{e}\left(  s\right)
\]


Sustituimos $i_{1}\left(  s\right)  :$%

\[
R_{1}[LCs^{2}+C(R_{2}+R_{3})s+1]i_{2}\left(  s\right)  +[L_{s}+\left(
R_{2}+R_{3}\right)  ]i_{2}\left(  s\right)  =V_{e}\left(  s\right)
\]


Factorizamos $i_{2}\left(  s\right)  :$%

\[
I_{2}(s)\{R_{1}[LCs^{2}+R_{1}C(R_{2}+R_{3})s+1]+L_{s}+\left(  R_{2}%
+R_{3}\right)  \}=V_{e}\left(  s\right)
\]


Desarollamos:%

\[
I_{2}(s)[R_{1}LCs^{2}+R_{1}C(R_{2}+R_{3})s+R_{1}+L_{s}+R_{2}+R_{3}%
=V_{e}\left(  s\right)
\]


Agrupamos por potencias de $s:$%

\[
I_{2}(s)[R_{1}LCs^{2}+(L+R_{1}C(R_{2}+R_{3}))s+(R_{1}+R_{2}+R_{3}%
)]=V_{e}\left(  s\right)
\]


Se obtiene $\frac{I_{2}(s)}{V_{e}\left(  s\right)  }:$%

\[
\frac{I_{2}(s)}{V_{e}\left(  s\right)  }=\frac{1}{R_{1}LCs^{2}+(L+R_{1}%
C(R_{2}+R_{3}))s+(R_{1}+R_{2}+R_{3})}
\]


Como la salida se encuentra en $R_{3}$ la salida es:%

\[
F_{s}(s)=R_{3}I_{2}\left(  s\right)
\]


Entonces...%

\[
\frac{F_{s}(s)}{V_{e}(s)}=R_{3}\frac{I_{2}(s)}{V_{e}(s)}
\]


Por lo tanto la funcion de transferencia final es de:%

\[
\frac{F_{s}(s)}{V_{e}(s)}=\frac{R_{3}}{R_{1}LCs^{2}+[L+R_{1}C(R_{2}%
+R_{3})]s+(R_{1}+R_{2}+R_{3})}
\]


\bigskip

\emph{ERROR\ EN\ ESTADO\ ESTACIONARIO.}

El error en estado estacionario es:%

\[
e_{ss}=1-\frac{R_{3}}{R_{1}+R_{2}+R_{3}}
\]
%

\[
e_{ss}=\frac{R_{1}+R_{2}+R_{3}-R_{3}}{R_{1}+R_{2}+R_{3}}
\]
%

\[
e_{ss}=\frac{R_{1}+R_{2}}{R_{1}+R_{2}+R_{3}}
\]


\emph{CASO.}

Sustituyendo:%

\[
e_{ss}=\frac{12+4.7}{12+4.7+18}\rightarrow e_{ss}=\frac{16.7}{34.7}=0.48
\]


\emph{CONTROL.}

Sustituyendo:%

\[
e_{ss}=\frac{1+1}{1+1+10}\rightarrow e_{ss}=\frac{2}{12}=0.167
\]


\emph{ESTABILIDAD DEL SISTEMA.}

\textbf{Control; Estable con respuesta sobreamortiguada.}%

\[
\lambda_{1}=-23.2116
\]
%

\[
\lambda_{2}=-54998.065
\]


\textbf{Caso; Estable con respuesta sobreamortiguada}%

\[
\lambda_{1}=-580057
\]
%

\[
\lambda_{2}=-11329.8731
\]



\end{document}